% ============================================================
%  chapters/ch01_balakanda.tex
%
%  LAYER STRUCTURE (mirrors template four-commentary layout,
%  reduced to one commentary layer per shloka for this work)
%  ─────────────────────────────────────────────────────────
%  \roottext{}              → shloka verse  (violet Devanagari)
%  \begin{commentary}{…}   → commentary heading (centred)
%  \end{commentary}
%  body text                → English / Hindi commentary
%
%  Add further \roottext + commentary blocks for new shlokas.
% ============================================================

% ── Activate page numbering ───────────────────────────────────
\fancyhead[LE]{\thepage}
\setcounter{page}{1}
\renewcommand{\headrulewidth}{1pt}

% ── Running headers ──────────────────────────────────────────
\fancyhead[CE]{\en\itshape Vālmīki Rāmāyaṇa~|}   % left page centre  ← change per chapter
\fancyhead[CO]{\en\itshape Bālakāṇḍa~|}           % right page centre ← change per chapter
\fancyhead[RO,LE]{\thepage}

% Fix for ch01_balakanda.tex
%
% The chapter-opening block had  \en नमः  which passes Devanagari text
% to Linux Libertine O (a Latin-only font), producing nullfont warnings.
% \en should only wrap Latin/IAST text.
%
% ── BEFORE (broken) ──────────────────────────────────────────────────────────
%
%   ~॥~\om\ \en नमः~॥\\
%
% ── AFTER (fixed) ────────────────────────────────────────────────────────────
%
%   \h{~॥~}\om\ \h{नमः~॥}\\
%
% \h{} is defined in preamble.tex as  \newcommand{\h}[1]{{\la #1}}
% so it selects Shobhika for its argument and returns to the default font.
% The standalone ॥ before \om was also unguarded; wrapping it in \h{} fixes
% that too.
%
% ── Full corrected chapter-opening block ─────────────────────────────────────

\begin{center}
{\Huge $_{*}$}~{\la श्रीः}~{\Huge $_{*}$}\\
\h{~॥~}\om\ \h{नमः~॥}\\
\vspace{2mm}
{\en\LARGE Bālakāṇḍa}\\
\rule{.1\linewidth}{1pt}\\
\vspace{1mm}
{\en\large The Book of Beginnings}
\end{center}

% ═══════════════════════════════════════════════════════════
%  SHLOKA 1
% ═══════════════════════════════════════════════════════════

\sarga{Sarga 1 · The First Shloka}

% ── Root text (shloka) ───────────────────────────────────────
\roottext{%
\om\ तपःस्वाध्यायनिरतं तपस्वी वाग्विदां वरम् ।\\
नारदं परिपप्रच्छ वाल्मीकिर्मुनिपुङ्गवम् ॥
}

% ── Commentary ───────────────────────────────────────────────
\begin{commentary}{\en Commentary}
\end{commentary}



\newpage
